\section{Introduction}\label{sec:introduction}
Write $\T=\R/\Z$ and $\e(t)=\exp(2\pi i t)$. All tori carry probability
Haar measure. A Littlewood polynomial of length $N$ is a polynomial
\[
 P(z)=\sum_{k=0}^{N-1}\varepsilon_k z^k,
 \qquad \varepsilon_k\in\{-1,1\}.
\]
Its maximum modulus is $\norm{P}_\infty=\max_{t\in\T}|P(\e(t))|$.
Define
\[
 m_N=\min_{\varepsilon_0,\ldots,\varepsilon_{N-1}\in\{-1,1\}}
       \frac{1}{\sqrt N}\norm{\sum_{k=0}^{N-1}\varepsilon_k z^k}_\infty.
\]
Orthogonality gives $\norm{P}_2^2=N$, so $m_N\ge1$. The real-sign
analogue of Erd\H{o}s's fixed-gap question asks whether $m_N$ is
eventually bounded below by $1+c$ for some absolute $c>0$.
Erd\H{o}s's original question allowed complex unimodular coefficients
and was refuted in that class by Kahane; the restriction to real signs
is the question addressed here \cite[Problem~22]{Erdos1957}; see also
\cite{Kahane1980} and \cite[Problems~4.13 and~4.31]{HaymanLingham2018}.

\begin{theorem}\label{thm:main}
For Littlewood polynomials with real coefficients,
\[
 \lim_{N\to\infty}m_N=1.
\]
Equivalently, for every $\eta>0$ there is $N_0$ such that for every integer
$N\ge N_0$ there are signs $\varepsilon_0,\ldots,\varepsilon_{N-1}$ with
\[
 \max_{|z|=1}\left|\sum_{k=0}^{N-1}\varepsilon_kz^k\right|
 \le(1+\eta)\sqrt N.
\]
\end{theorem}

The theorem concerns every sufficiently large length, rather than a
subsequence. It gives no uniform lower bound for $|P(z)|$ and hence does
not assert two-sided uniform ultraflatness. The choices in the proof are
existential; no useful convergence rate or efficient signing algorithm
is asserted.

\subsection{Context and comparison}
The Shapiro--Rudin construction, obtained earlier by Shapiro and
independently by Rudin, gives real-sign polynomials
with maximum $O(\sqrt N)$; at dyadic lengths its paired identity gives
an upper bound $\sqrt{2N}$ \cite[Theorem~I and its proof]{Rudin1959}.
Balister proved the
all-length bound $\sqrt{6N-2}-1$ for its initial segments
\cite[Theorem 1]{Balister2019}. A different problem asks for upper and
lower bounds by fixed positive multiples of $\sqrt N$ throughout the
circle. Such two-sided flat Littlewood polynomials were constructed by
Balister, Bollob\'as, Morris, Sahasrabudhe and Tiba, whose proof combines
the Shapiro--Rudin construction with discrepancy theory
\cite{BalisterEtAl2020}. These results control the scale of the maximum,
but do not make the upper constant tend to one.

The coefficient class matters. Complex coefficients of modulus one allow
ultraflat constructions. Littlewood used quadratic phases near this
Parseval scale, Kahane obtained ultraflatness, and Bombieri and Bourgain
gave quantitative constructions \cite{Littlewood1966,Kahane1980,BombieriBourgain2009}.
The smoothed quadratic phases and Poisson summation used in
\cite[Sections~2 and~7]{BombieriBourgain2009} are analytic predecessors
of the sampling argument below. Their bounded-to-unimodular correction
also develops a method of K\"orner
\cite[Section~8]{BombieriBourgain2009}. The present proof must instead
preserve real coefficients and round them to signs. The distinction
between the unimodular and real-sign questions is also explicit in
\cite[Problem 4.31]{HaymanLingham2018}.
Georgiev, G\'omez-Serrano, Tao and Wagner recently used AlphaEvolve to
search for the same real-sign minimum at degrees up to $100$;
their computations do not give an asymptotic theorem
\cite[Section 6.13]{GeorgievEtAl2025}.

There can still be a growing additive gap above Parseval's bound.
Erd\'elyi proved that every Littlewood polynomial of length $N$ satisfies
\cite[Theorem 2.1]{Erdelyi2026}
\[
 \norm{P}_\infty^2\ge N+\frac{(N-1)^{1/3}}{38}.
\]
The extra term divided by $N$ tends to zero, so this is compatible with
Theorem~\ref{thm:main}.

The theorem implies flatness in every fixed finite $L^p$ modulus
sense. This conflicts with nonflatness claims in preprints of
el Abdalaoui \cite{ElAbdalaoui2017,ElAbdalaoui2025,ElAbdalaoui2025Generalization}.
Section~\ref{sec:finite-flatness} proves the implication with our
normalization. Appendix~\ref{sec:concentration} examines specific issues in
the contrary arguments.

\subsection{Merit factors and a symbolic spectral consequence}
For a binary word $A=(a_0,\ldots,a_{N-1})\in\{-1,1\}^N$ of length
$N\ge2$, define its aperiodic sidelobe autocorrelations and merit factor by
\[
 C_u(A)=\sum_{j=0}^{N-1-u}a_ja_{j+u}\quad(1\le u<N),
 \qquad
 F(A)=\frac{N^2}{2\sum_{u=1}^{N-1}C_u(A)^2}.
\]
This is the standard merit-factor normalization
\cite[Definition~2 and Lemma~0]{DownarowiczLacroix1998}; see also
\cite[Introduction]{JedwabKatzSchmidt2013} for its fourth-moment formulation.
The conjectured boundedness of $F(A)$ over all such words is called
Turyn's conjecture, equivalently Erd\H{o}s's $L^4$-norm conjecture, in
\cite[p.~3 of the author manuscript]{DownarowiczLacroix1998}.
For comparison, Jedwab, Katz and Schmidt constructed families with
limiting merit factor $6.342061\ldots$, improving the
previous proven value $6$, by modifying and extending Legendre sequences
\cite[Theorem~1.1 and Corollary~3.2]{JedwabKatzSchmidt2013}.
Section~\ref{sec:merit-morse} disproves this boundedness in the stronger
all-length form: the largest merit factor at length $N$ tends to infinity
as $N$ tends to infinity through all integers.

Downarowicz and Lacroix also relate this question to binary Morse shifts,
the discrete-time symbolic systems generated by iterated products of binary
words. Their theorem converts unbounded binary merit factors into a uniquely
ergodic binary Morse shift with simple spectrum whose zero-coordinate
spectral measure is absolutely continuous with an $L^2$ density
\cite[Theorem~2 and its proof]{DownarowiczLacroix1998}.
Section~\ref{sec:merit-morse} gives the precise spectral statement and its
deduction from Theorem~\ref{thm:main}.

\subsection{Proof overview}
We first allow real coefficients $X_{N,k}\in[-1,1]$. Their defect from
being signs is measured by
\[
 \mu(X_N)=\frac12\sum_{k=0}^{N-1}(1-|X_{N,k}|).
\]
The construction makes the normalized Fourier maximum close to one
while keeping $N^{-1}\sum_k X_{N,k}^2$ close to one. The latter condition
makes $\mu(X_N)/N$ small. A discrepancy argument then replaces the
coefficients by signs, with an error controlled by this small defect.
Section~\ref{sec:reduction} states both intermediate results precisely
and gives the deduction of the theorem.

To build the relaxed coefficients, we use a real trigonometric polynomial
$F$ on an auxiliary torus whose absolute value is at most one and whose
mean-square mass is nearly one. Each Fourier mode of $F$ is assigned a
nonzero quadratic curvature. A stationary-phase calculation shows that
sampling such a mode on a block spreads its Fourier contribution over
an interval whose width is proportional to that curvature. There are two
requirements: distinct modes and blocks must have disjoint intervals,
and each individual contribution must have size at most approximately
$\sqrt N$.

Section~\ref{sec:packing} supplies the interval packing, even when the
integer frequency vectors have linear relations. It uses a randomly
constructed finite-field hypergraph and the Pippenger--Spencer coloring
theorem, in the near-perfect-packing lineage of Frankl and R\"odl
\cite{FranklRodl1985,PippengerSpencer1989,AlonYuster2005}.
Section~\ref{sec:spreading} adjusts the Fourier coefficients of $F$ by
quadratic oscillations in extra variables so that their sizes match the
available widths. Section~\ref{sec:sampling} samples the resulting fixed
polynomial on disjoint blocks. At each angle at most one stationary
contribution remains, whereas orthogonality in the sampling limit
recovers nearly full mean-square mass.

Finally, Section~\ref{sec:rounding} proves the defect-sensitive rounding
bound. We use the partial-coloring method of Beck and Spencer
\cite{Beck1981,Spencer1985}, in the arbitrary-vector form verified through
Lovett--Meka's theorem \cite{LovettMeka2015}. A dyadic rounding procedure
keeps the defect mass from increasing, which makes its total loss
$o(\sqrt N)$ as the target error decreases. All auxiliary dimensions,
frequencies, primes, blocks and smooth cutoffs are fixed before $N$
tends to infinity.
